%----------------------------------------------------------------------------------------
%    Solution
%----------------------------------------------------------------------------------------
\section{UML Diagrams}
\subsection{Activity Diagrams:}
\begin{figure}[h]
    \centering
    \includegraphics[width=1\textwidth]{ActivityDiagram.png}
    \label{fig:ActivityDiagrams}
\end{figure}
\noindent These diagrams represent the optimization algorithm. The process happens inside the optimizer class, it starts when the GetResultData method of the optimizer is called. That calls the Optimize method, which proceeds to retrieve all data necessary for optimization to commence. 

\noindent When data has been gathered it enters a loop, going through every hour of every day of the optimization period. In the loop it checks if the gathered data is workable, after, if data is correct, the process continues. The weighted ratio of heat to price of each unit is calculated, then the optimal unit is selected to be used in heat production and usage is recorded. This happens on a per hour basis by calling CalculateHeatPerPrice and ChooseAndRecordUsage methods.

\noindent CalculateHeatPerPrice method loops through the units, for each unit the ratio is calculated by first getting a crude heat divided by price, then the crude ratio is weighted by a factor of ten if it is negative. The weighting is performed to increase the unit's priority, because negative heat per price means that the unit is making money instead of costing money. 

\noindent The ChooseAndRecordUsage method first sorts units by their heat per price ratio. When the list is sorted it enters a loop going through the sorted list, keeping track of heat demand, how much heat has been produced and how many units are available. Loop stops if it runs out of units or if the heat demand is reached. At the start of the loop it checks if the current unit is on maintenance, if unit is on maintenance it is skipped, otherwise loop proceeds. The next step is to check if the unit's maximum heat production exceeds remaining heat demand, which is the difference of total heat demand and produced heat. This is done to calculate workload - a coefficient from 0 to 1 that represents how much unit is used. If heat production is not exceeded the workload is 1, otherwise it is calculated by dividing remaining heat demand by the current unit's maximum heat production. After calculating workload, produced heat is increased by current unit’s maximum heat production multiplied by workload, and then all unit’s relevant characteristics, such as pollution, heat, fuel consumption, are recorded multiplied by workload.

\subsection{Sequence Diagram:}
\begin{figure}[h]
    \centering
    \includegraphics[width=1\textwidth]{SequenceDiagram.png}
    \label{fig:SequenceDiagram}
\end{figure}
\noindent This sequence diagram represents how service classes interact when result data is requested.

\noindent Requests can come from several parts of the Presentation layer, which is represented as a found message.

\noindent The Result Data Manager makes a self-call, which then calls Optimizer. After that is resolved, within the same self-call another self-call is made to save Optimization results in local .csv files. For data management a call to a separate class is made, represented as a lost message. After that, the Result Data Manager’s sequence ends. 

\noindent The Optimizer makes a self-call that then makes three self-calls, first resulting in a call to the Asset Manager and second in two calls to the Source Data Manager. It then enters a loop of two self-calls. After that Optimizer’s sequence ends.

\noindent The Asset Manager makes a self-call to initialize. Initialization uses a separate class for data management, this is represented as a lost message. After that is resolved Asset Manager uses two separate classes to provide the correct scenario, and units of correct classes - those are represented as lost messages. After that Asset Manager’s sequence ends.

\noindent Source Data Manager sequence starts when it receives the first call from Optimizer. This results in a self-call that uses a separate class for data manipulation represented as two lost messages. Then a second call from the optimizer results in following the same sequence. After that Source Data Manager’s sequence ends.

\subsection{Class Diagram:}
\begin{figure}[h]
    \centering
    \includegraphics[width=1\textwidth]{ClassDiagram.png}
    \label{fig:ClassDiagram}
\end{figure}
\noindent This class diagram represents the core of the solution: services and models.

\noindent AssetManager has a one-way association with the Optimizer. This one-way association shows that AssetManager is called by the Optimizer, but doesn’t call it. AssetManager aggregates instances of GenericUnit and Grid. This aggregation relationship shows that those are independent concepts.

\noindent SourceDataManager has a one-way association with the Optimizer. This shows that Optimizer can call SourceDataManger, but SourceDataManager doesn’t call Optimizer. It aggregates the TimeSeries<T> class. This represents that TimeSeries<T> is an independent class.

\noindent Optimizer has one-way associations with AssetManager, SourceDataManager and ResultDataManager. Those represent that Optimizer calls AssetManager and SourceDataManager, but is not called by them. And that shows that Optimizer is called by ResultDataManager, but doesn’t call it. Optimizer aggregates TimeSeries<T> and OptimizedUnit. The aggregation relationship shows that both are independent concepts.

\noindent ResultDataManager has a one way association with the Optimizer. This shows that Optimizer is called by ResultDataManager, but ResultDataManager is not called by Optimizer. ResultDataManager  aggregates OptimizedUnit. This aggregation relationship shows that this is an independent class.

\noindent Unit models form an inheritance tree. GenericUnit class and abstract OptimizedUnit classes inherit from the base abstract class Unit. GenericUnit represents units with characteristics including an image, while OptimizedUnit represents units with only performance specifications, and records of usage. From the abstract OptimizedUnit class three classes: CombustionUnit, HybridUnit and ElectricUnit, inherit to allow for the differences in how they calculate price and how they record usage.

\section{Technical Architecture}
\noindent Project was coded in C\# in .NET framework, using Avalonia UI in combination with OxyPlot, ReactiveUI and CommunityToolkit.Mvvm for Data Visualization. All input and output data, specifically: images, units' specifications and resulting schedule are stored locally, file manipulation is performed by use of System library and CsvHelper library. Tests were implemented in the xUnit framework.

\noindent Solution is separated in following layers:
\begin{itemize}

    \item [\(\cdot \)] Data Handling
    \begin{itemize}
        \item [\(\circ \)] Data
        \item [\(\circ \)] Assets
    \end{itemize}
    \item [\(\cdot \)] Business Logic
    \begin{itemize}
        \item [\(\circ \)] Models
        \item [\(\circ \)] Services
        \item [\(\circ \)] ViewModels
    \end{itemize}
    \item [\(\cdot \)] Presentation
    \begin{itemize}
        \item [\(\circ \)] Views
    \end{itemize}

\end{itemize}
\noindent The project's architecture is a combination of the MVVM pattern and micro-service architecture approach.

\noindent Data Handling layer includes: Data - classes that handle file manipulation; Assets - static data storage consisting of .json, .csv and .png files.

\noindent Business Logic includes: Models - classes that: model real life objects and concepts, transfer data, represent UI components; Services - classes that manage classes from Data, and control the process of optimization; ViewModels - classes that provide functionality to the UI and establish connection between Models, Services and Presentation layer.

\noindent Presentation layer includes: Views - partial classes that control the appearance of the UI.

\noindent This structure was chosen because it follows SOLID principles which facilitate flexibility and scalability, by significantly reducing required amount of refactoring in development process. MVVM pattern was implemented in the overall structure of this project. Core benefit of MVVM is separation of responsibilities. Microservice architecture is most represented by Services and it follows the Open/Closed principle.
